\chapter{Vision General del Dataset}
\label{cap:overview}

\section{Caracteristicas Principales}

El dataset FatigueSet contiene datos fisiologicos y comportamentales de
\textbf{12 participantes} que realizaron un protocolo experimental estructurado
para inducir fatiga mental y fisica controlada. Los datos fueron recopilados
en \textbf{3 sesiones} por participante, y cada sesion se dividio en
\textbf{3 fases de medicion}.

Las dimensiones principales del dataset son:

\begin{table}[H]
\centering
\caption{Dimensiones principales del FatigueSet}
\begin{tabular}{|l|r|p{5cm}|}
\hline
\textbf{Dimension} & \textbf{Valor} & \textbf{Notas} \\
\hline
Participantes & 12 & Etiquetados P01 a P12 \\
\hline
Sesiones por participante & 3 & Sesiones 01, 02, 03 \\
\hline
Fases de medicion por sesion & 3 & M1, M2, M3 \\
\hline
Dispositivos simultaneos & 4 & Wearables (ver Cap. 3) \\
\hline
Tareas cognitivas & 5 & CRT, N-Back, Task Switch, Fatigue Scale, Markers \\
\hline
Variables fisiologicas agregadas & 19 & HR, BR, HRV, EDA, EEG, etc. \\
\hline
Reglas de combinacion total & 12 x 3 x 3 & 108 registros en archivo agregado \\
\hline
\end{tabular}
\end{table}

\section{Protocolo Experimental y Contrabalanceo}

\subsection{Intensidad de Actividad}

Cada sesion fue disenada para inducir un nivel diferente de fatiga fisica:

\begin{itemize}
    \item \textbf{Low Intensity}: Actividad fisica ligera (caminar lentamente o
          ejercicio estatico suave). Induce minimo esfuerzo y fatiga.

    \item \textbf{Medium Intensity}: Actividad fisica moderada (ejercicio cardiovascular
          moderado). Induce fatiga moderada pero sostenible.

    \item \textbf{High Intensity}: Actividad fisica vigorosa (ejercicio intenso o HIIT-style).
          Induce fatiga aguda pronunciada.
\end{itemize}

\subsection{Contrabalanceo de Orden}

Para evitar efectos del orden (aprendizaje, fatiga acumulada entre sesiones),
los participantes fueron asignados a \textbf{ordenes contrabalanceados} de intensidad.
Por ejemplo:

\begin{table}[H]
\centering
\caption{Ejemplo de contrabalanceo (P01)}
\begin{tabular}{|c|c|c|c|}
\hline
\textbf{Participante} & \textbf{Sesion 1} & \textbf{Sesion 2} & \textbf{Sesion 3} \\
\hline
P01 & Low (1) & Medium (2) & High (3) \\
P02 & Low (1) & High (3) & Medium (2) \\
P03 & Medium (2) & High (3) & Low (1) \\
P04 & High (3) & Medium (2) & Low (1) \\
\vdots & \vdots & \vdots & \vdots \\
\hline
\end{tabular}
\end{table}

Este contrabalanceo de cuadrado latino asegura que no hay confusiones sistematicas
entre orden de sesion e intensidad de actividad.

\section{Tres Fases de Medicion: M1, M2, M3}

Dentro de cada sesion, el protocolo incluyo \textbf{tres fases temporales} donde
se recopilaron mediciones de todos los sensores:

\begin{description}
    \item[M1 - Baseline (Linea Base)]
    Duracion: ~10 minutos. El participante se encuentra en reposo, se sincroniza
    la instrumentacion, se recopilan mediciones EEG basales (ojos abiertos y cerrados).
    Esta fase sirve de comparacion para detectar cambios inducidos por la actividad.

    \item[M2 - Post-Ejercicio]
    Duracion: ~5-20 minutos (variable segun la intensidad de la actividad previa).
    Inmediatamente despues de completar la actividad fisica designada. Es la fase
    donde se observan cambios mas marcados en variables cardiovasculares y electrodermal.

    \item[M3 - Post-Fatiga Mental]
    Duracion: ~10 minutos. Despues de un bloque de tareas cognitivas (CRT, N-Back, Task Switch).
    Permite evaluar fatiga mental e interacciones entre fatiga fisica y mental.
\end{description}

Durante cada fase, los participantes completaron:
\begin{itemize}
    \item \textbf{Escala de Fatiga Subjetiva} (exp\_fatigue.csv) - 2 preguntas (fisica y mental)
    \item \textbf{Tareas Cognitivas} (exp\_crt.csv, exp\_nback.csv, exp\_task\_switch.csv)
    \item \textbf{Recopilacion de Datos Fisiologicos} - 4 dispositivos simultaneamente
\end{itemize}

\section{Estructura de Archivos}

El dataset esta organizado jerarquicamente:

\begin{verbatim}
fatigueset/
├── metadata.csv                          (Informacion de contrabalanceo)
├── 01/ 02/ ... 12/                       (Directorios por participante)
│   ├── 01/ 02/ 03/                       (Subdirectorios por sesion)
│   │   ├── exp_crt.csv                   (Tarea CRT)
│   │   ├── exp_nback.csv                 (Tarea N-Back)
│   │   ├── exp_task_switch.csv           (Tarea Task Switch)
│   │   ├── exp_fatigue.csv               (Escala de fatiga subjetiva)
│   │   ├── exp_markers.csv               (Marcadores temporales)
│   │   ├── ear_acc_left.csv, ...         (Nokia eSense - 6 archivos)
│   │   ├── forehead_eeg_alpha_abs.csv, ...(Muse S - 13 archivos)
│   │   ├── chest_physiology_summary.csv, (Zephyr BioHarness - 9 archivos)
│   │   └── wrist_eda.csv, ...            (Empatica E4 - 6 archivos)
└── README.md                             (Descripcion tecnica completa)

Total: 34 archivos CSV por sesion
\end{verbatim}

\section{Dataset Agregado: Caracteristicas Derivadas}

Aunque los archivos crudos contienen miles de registros por sesion (en frecuencias
que van de 1 Hz a 256 Hz), se proporciona un archivo agregado a nivel de sesion:

\begin{verbatim}
fatigueset_aggregated_features.csv
\end{verbatim}

Este archivo contiene \textbf{19 columnas} con estadisticas resumidas (media y desviacion
estandar) de cada variable fisiologica por registros por fase y sesion:

\begin{table}[H]
\centering
\caption{Estructura del archivo agregado (19 columnas)}
\begin{tabular}{|l|l|l|}
\hline
\textbf{Columnas de Identificacion} & \textbf{Autoreportes} & \textbf{Metricas Fisiologicas} \\
\hline
participante & fatiga\_fisica & hr\_media, hr\_std \\
sesion & fatiga\_mental & br\_media, br\_std \\
intensidad & & hrv\_media, hrv\_std \\
intensidad\_num & & eda\_media, eda\_std \\
fase & & eeg\_alpha\_media \\
fase\_num & & eeg\_beta\_media \\
& & eeg\_theta\_media \\
\hline
\end{tabular}
\end{table}

\textbf{Forma del dataset agregado}: 108 filas (12 participantes x 3 sesiones x 3 fases)

\section{Resumen Estadistico}

Al limpiar el dataset de valores inf/NaN, disponemos de \textbf{104 registros validos}
cuando se requieren todas las variables simultaneamente.

\begin{table}[H]
\centering
\caption{Estadisticas descriptivas globales}
\begin{tabular}{|l|l|l|l|}
\hline
\textbf{Variable} & \textbf{Media} & \textbf{Desv. Std.} & \textbf{Rango} \\
\hline
HR (bpm) & 88.35 & 20.35 & [58.91, 162.28] \\
BR (brpm) & 20.87 & 4.56 & [11.68, 33.43] \\
HRV (ms) & 69.23 & 31.18 & [11.28, 189.38] \\
EDA (µS) & 1.30 & 1.82 & [0.08, 7.68] \\
EEG Alpha & 0.575 & 0.159 & [0.039, 0.998] \\
EEG Beta & 0.402 & 0.144 & [0.030, 0.753] \\
EEG Theta & 0.436 & 0.198 & [0.027, 1.054] \\
\hline
\end{tabular}
\end{table}

Estos rangos y promedios seran interpretados en detalle en los capitulos posteriores.
